% ============================================================================
% sections/system.tex -- Section 3: System and Failure Model
% Drafted from the P1-approved docs/method.md (§1.3--§1.5, §2, §3).
% ============================================================================

% Theorem-like environments. Defined once here (guarded, so a later preamble
% addition cannot clash); base-LaTeX \newtheorem is used because amsthm is not
% loaded by the class.
\makeatletter
\@ifundefined{c@proposition}{\newtheorem{proposition}{Proposition}}{}
\@ifundefined{c@lemma}{\newtheorem{lemma}{Lemma}}{}
\@ifundefined{c@corollary}{\newtheorem{corollary}{Corollary}}{}
\@ifundefined{proofsketch}{%
  \newenvironment{proofsketch}{\par\noindent\emph{Proof sketch.}\ }{\hfill$\square$\par}}{}
\makeatother

\section{System and Failure Model}
\label{sec:model}

This section fixes the vocabulary used throughout the paper: the pipeline and its
durable stores, the event and ledger schemas, the ledger invariants, the boundary of
Kafka's exactly-once guarantee, and the crash-point model against which the mechanisms
are analysed.

\subsection{Pipeline and durable stores}
\label{sec:model:pipeline}

The pipeline is the canonical form of a ledger ingestion service: a Kafka source topic
carries business events; a consumer service applies a deterministic function to each
record and performs two durable side effects; a load generator produces the events and
an independent oracle verifies the end state. Kafka is used as a replicated,
partitioned log \cite{kafka2015}; the source topic is partitioned, replicated with
factor three and consumed by a consumer group whose offsets are committed to Kafka's
offset store. The consumer group is the unit of progress; partitions are the unit of
parallelism.

Two durable stores receive side effects. The first is the \emph{ledger sink}: a
PostgreSQL database whose posting table is the authoritative record of the ledger. The
second is the \emph{publication stream}: a Kafka topic to which one derived entry record
per business event is appended, used by downstream audit consumers
\cite{sax2018,eventsourcing2022}. For the DB-sink mechanisms (M0, M1x, M2, M3, M4 and
the combinations) the ledger is PostgreSQL and the publication stream is auxiliary; for
the Kafka-sink mechanism (M1-kafka) the publication stream \emph{is} the ledger. This
distinction is a construct-validity matter and is stated wherever mechanisms are
compared (Section~\ref{sec:mechanisms}).

The deployment analysed here is deliberately modest and reproducible: three Kafka
KRaft broker processes, one PostgreSQL instance, one Toxiproxy process for TCP faults,
and the consumer service at two instances in one consumer group. Table~\ref{tab:components}
lists the components and their failure-relevant roles. The single-host deployment is a
fidelity limitation discussed in the threats-to-validity section; it is not a modelling
simplification: the mechanisms are the ones deployed in practice, and every fault is a
real process kill or a real TCP manipulation, not a simulation.

\begin{table*}[t]
\centering
\small
\caption{Executed system: components, roles, and failure-relevant properties. All
components are native processes; the substrate and version pinning are described in
Section~\ref{sec:method}.}
\label{tab:components}
\begin{tabular}{@{}p{2.6cm}p{2.2cm}p{6.2cm}p{4.2cm}@{}}
\toprule
Component & Instances & Role & Durable state \\
\midrule
Kafka brokers (KRaft) & 3 & replicated source log, publication stream, consumer offsets, transaction state & log segments, offset store, transaction state \\
PostgreSQL 16 & 1 & authoritative ledger for DB-sink mechanisms; progress, inbox and outbox tables where the mechanism uses them & ledger tables, progress tables, WAL \\
Consumer service & 2 & one instance of the mechanism under test; reads source records and performs the mechanism's durable duties & none of its own (recovery uses the mechanism's durable state) \\
Load generator & 1 & deterministic open-loop producer; writes ground truth before producing & ground-truth file and produced manifest \\
Toxiproxy & 1 & cuts, delays and resets real TCP connections (app$\leftrightarrow$Kafka, app$\leftrightarrow$PostgreSQL) & none \\
Runner and injectors & 1 & per-run orchestration, fault injection, fault attestation, run records & run records and fault log \\
Oracle & 1 & independent verifier: re-derives the expected postings and compares them with the sink & analysis output only \\
\bottomrule
\end{tabular}
\end{table*}

\subsection{Events, entries, and postings}
\label{sec:model:schema}

One source record carries exactly one complete journal entry: a business event
$e=(\mathit{eid},\mathit{cur},\mathit{legs})$ with $\mathit{legs}=((a_D,D,m),(a_C,C,m))$,
$m\in\mathbb{Z}^{+}$ and $a_D\neq a_C$. The two legs are one debit (D) of $m$ minor
units and one credit (C) of the same amount in the same currency. The restriction to
two-leg, single-currency entries is a scoping decision: it keeps the unit of atomicity
at the record level and isolates the external-store boundary, which is the object of
study; multi-leg and multi-currency entries are not modelled
(Section~\ref{sec:model:assumptions}, assumption A12).

Events are a pure function of a per-repetition seed and a global sequence number: the
$\mathit{eid}$, both account identifiers, the amount and the target partition are
derived with SHA-256 over domain-separated inputs. The generator and the oracle
implement the same derivation independently, so the expected posting set is
recomputable rather than stored, and the two implementations can be compared
byte-for-byte. This is the mechanism behind the independent oracle of
Section~\ref{sec:method}.

The ledger sink stores \emph{postings}, one row per leg. Table~\ref{tab:posting}
summarises the schema. Postings are append-only: the application role holds only
\texttt{INSERT} and \texttt{SELECT} on the table, and a trigger rejects every
\texttt{UPDATE} or \texttt{DELETE}. The deduplication index on
$(\mathit{event\_id},\mathit{leg\_no})$ exists only in mechanisms that use
de-duplication; its absence in the naive mechanisms is part of the treatment, not an
implementation accident. Amounts are integers in minor units; no floating point
appears in the pipeline.

\begin{table}[t]
\centering
\small
\caption{Ledger sink schema (columns abbreviated). The unique index is present only in
the de-duplicating mechanisms (M3, M2+M3, M4+M3, M1x+M3). Append-only enforcement is
identical for all mechanisms.}
\label{tab:posting}
\begin{tabular}{@{}p{2.9cm}p{1.8cm}p{2.4cm}@{}}
\toprule
Column & Type & Constraint \\
\midrule
\texttt{posting\_id} & bigserial & primary key \\
\texttt{event\_id} & text & not null \\
\texttt{leg\_no} & smallint & check $\in\{1,2\}$ \\
\texttt{account\_id} & text & not null \\
\texttt{side} & char(1) & check $\in\{\mathrm{D},\mathrm{C}\}$ \\
\texttt{amount\_minor} & bigint & check $>0$ \\
\texttt{currency} & char(3) & not null \\
\texttt{produced\_at} & timestamptz & from the event \\
\midrule
\multicolumn{3}{@{}p{7.1cm}@{}}{\texttt{UNIQUE(event\_id, leg\_no)} --- de-duplicating mechanisms only} \\
\bottomrule
\end{tabular}
\end{table}

\subsection{Ledger invariants and observation}
\label{sec:model:invariants}

A ledger is correct when it records every consumed business event exactly once and no
other posting. Double-entry bookkeeping requires each entry to be balanced
\cite{triple2019,triple2021}; append-only recording is the classic durability
discipline for journals \cite{lsfs1992}. We therefore identify four invariants. Let
$D$ be the observed posting multiset (rows of the sink, or committed records of the
publication stream when that stream is the ledger) and let $\Omega$ be the expected
posting multiset derived from the events that are durably present in the source topic.
Write $\mathrm{mult}(D,\pi)$ for the multiplicity of posting $\pi=(e,\ell)$ in $D$.
Table~\ref{tab:invariants} states the four invariants and their executable form.

I1 is split into structural well-formedness (\emph{I1a}: exactly two legs, one per
side, equal amount, single currency) and the signed-sum condition (\emph{I1b}); the
split matters because whole-entry duplication leaves the signed sum at zero. I4 is an
end-state completeness property, measured after quiescence; the corresponding liveness
argument is stated, but not relied upon for measurement.

\begin{table*}[t]
\centering
\small
\caption{The four ledger invariants and their executable form (abbreviated).
\texttt{expected} is the oracle's expected multiset $\Omega$; \texttt{signed} is
$+\mathit{amount}$ for a debit and $-\mathit{amount}$ for a credit. All predicates are
evaluated by the independent oracle (Section~\ref{sec:method}).}
\label{tab:invariants}
\begin{tabular}{@{}p{0.9cm}p{7.0cm}p{7.1cm}@{}}
\toprule
 & Formal statement & Executable predicate (summary) \\
\midrule
I1 & For every $e\in D$: exactly one D and one C leg, equal amounts, one currency,
signed sum zero. & group postings by event id; flag groups with $\mathit{count}\neq2$,
or $\sum\mathit{signed}\neq0$, or more than one currency, or not exactly one D and one
C. \\
I2 & No durable posting is modified or deleted, and every matched posting carries the
expected value. & compare $(\mathit{account},\mathit{side},\mathit{amount},
\mathit{currency})$ with the expected value on matched keys; attest that the application
role may not \texttt{UPDATE}/\texttt{DELETE} and that the append-only trigger exists. \\
I3 & $\forall e\in D:\ \mathrm{mult}(D,e)\le 2$ and $\forall\pi:\
\mathrm{mult}(D,\pi)\le 1$. & flag keys with count $>1$ per
$(\mathit{event\_id},\mathit{leg\_no})$, and events with more than two postings. \\
I4 & $\forall e\in E_{\mathrm{source}}$: the two postings of $\Pi(e)$ occur exactly
once in $D$. & anti-join of the expected multiset against the sink: every expected key
must be present exactly once. \\
\bottomrule
\end{tabular}
\end{table*}

The four invariants are deliberately orthogonal. A duplicated balanced entry is still
balanced: it violates I3, not I1. A missing entry is invisible to a uniqueness index: it
violates I4, not I3. A mutated amount may violate I2 while the entry remains
well-formed. This orthogonality is the analytical basis for the mechanisms' failure
signatures in Section~\ref{sec:formal}.

After the application is stopped and quiescence is established, the observation
function returns the sink multiset, the publication-stream multiset, the durable
progress state and the timing record; it is taken in one repeatable-read snapshot
\cite{pg-isolation,berenson1995}. A run is \emph{consistent} when
\begin{equation}
\label{eq:consistency}
\mathrm{Consistent}(D,\Omega)\ \Longleftrightarrow\ \text{I1}\wedge\text{I2}\wedge\text{I3}\wedge\text{I4}.
\end{equation}
The predicates are evaluated by the independent oracle (Section~\ref{sec:method});
their SQL realisation is released with the artifact.

\subsection{The exactly-once boundary}
\label{sec:model:eos}

Kafka's exactly-once semantics for a consume, transform and produce pipeline rest on
the transactional producer \cite{kip98,kip447,kafka-design}. A transaction may include
records appended to Kafka topics and the consumer offsets of the records being
processed, and the transaction coordinator makes both visible atomically. Read-committed
consumers observe all or none of the transaction, and stale producer epochs are fenced.
Formally, let $H_K$ be the durable sub-history at the Kafka site and let a transaction
$T$ contain the produced records and the processed offsets. The set of transitions that
Kafka makes atomic is then
\begin{multline}
\mathrm{EOS}(T)\ =\ \mathrm{vis}\big(\mathrm{Append}(L)\\
\cup\ \mathrm{CommitOffsets}(G)\big),
\label{eq:eos}
\end{multline}
a single-site atomic visibility relation inside Kafka.

The ledger, however, lives at a different site. Its commit
$\mathrm{Commit}_{S}$ is a PostgreSQL transaction that shares no coordinator, no
prepared state and no recovery protocol with Kafka. Atomic commit across independent
sites requires a distributed commit protocol with presumed-abort recovery
\cite{txncommit2006}; PostgreSQL can participate in externally coordinated prepared
transactions \cite{pg-prepare}, but the analysed pipeline does not, and the Kafka
transaction is not a two-phase-commit participant with the database. Compensating
transactions \cite{sagas1987} and application-level coordination
\cite{laigner2021,micro2026} are the alternatives practitioners use; they do not make
the two commits atomic either.

Consequently, for an event $e$ processed in a batch $b$, two windows exist. Let
$\Sigma$ be the set of admissible schedules, $\mathrm{Dur}_S(b)\subseteq\Sigma$ the
schedules in which the effect commit of $b$ is durable, and
$\mathrm{Dur}_K(b)\subseteq\Sigma$ the schedules in which the progress commit of $b$ is
durable. Then
\begin{align}
W_1(b) &= \mathrm{Dur}_S(b)\setminus\mathrm{Dur}_K(b), \label{eq:w1}\\
W_2(b) &= \mathrm{Dur}_K(b)\setminus\mathrm{Dur}_S(b). \label{eq:w2}
\end{align}
$W_1$ is the \emph{duplicate} window: the effect was applied but progress did not
advance, so a redelivery re-applies the effect unless the effect is idempotent. $W_2$
is the \emph{loss} window: progress advanced but the effect was not applied, so the
event is never redelivered and the posting is lost. Lemma~\ref{lem:windows} records
what a pair of independent single-site commits can and cannot rule out.

\begin{lemma}[Single-site commit pair]
\label{lem:windows}
Let a mechanism execute $\mathrm{Commit}_{S}$ and $\mathrm{Commit}_{K}$ as two
sequential durable operations with no cross-site recovery protocol. If the program
order is $\mathrm{Commit}_{S}\prec\mathrm{Commit}_{K}$, then $W_1\neq\emptyset$; if it
is $\mathrm{Commit}_{K}\prec\mathrm{Commit}_{S}$, then $W_2\neq\emptyset$. If the two
commits are concurrent or their order is not enforced, both windows may be non-empty.
\end{lemma}

\begin{proofsketch}
Under crash-recovery semantics with prefix-consistent storage
(assumption A1), a crash scheduled between two sequential durable operations leaves the
first durable and the second not. A process kill is such a crash. Concurrency admits an
interleaving in which either commit is durable alone.
\end{proofsketch}

\begin{corollary}[The EOS gap]
\label{cor:eosgap}
A pipeline that wraps its Kafka-internal appends and offsets in one Kafka transaction
but commits the external ledger with a plain database commit leaves $W_1$ open, and
$W_2$ open under weak or unenforced ordering. Kafka's atomicity does not cross the
site boundary of Eq.~\eqref{eq:eos}.
\end{corollary}

Figure~\ref{fig:eos} depicts the two windows on a timeline; the traces that realise
them are given in Section~\ref{sec:formal}. Kafka's guarantee does remain exactly-once
\emph{inside} the site: the records published to the publication stream are either all
visible or none, which is what makes a Kafka-sink mechanism a valid positive control
and an external-DB mechanism with a Kafka publication stream capable of \emph{diverging}
from its ledger.

\begin{figure}[t]
\centering
\resizebox{\linewidth}{!}{%
\begin{tikzpicture}[font=\footnotesize,>=Stealth,
  store/.style={draw,rounded corners,minimum height=0.8cm,align=center,inner sep=2pt},
  lab/.style={align=center}]
\draw[->] (0,0) -- (8.1,0) node[right]{time};
\node[store,anchor=west,minimum width=7.3cm] (k) at (0,1.45)
  {Kafka log + consumer offsets (site $K$)};
\node[store,anchor=west,minimum width=7.3cm] (p) at (0,-1.45)
  {PostgreSQL ledger (site $S$)};
\draw[->,thick] (2.9,1.0) -- (2.9,0.15) node[lab,above right=1pt,pos=0.2]{produce};
\draw[->,thick] (3.4,-1.0) -- (3.4,-0.15) node[lab,below right=1pt,pos=0.2]{DB COMMIT};
\draw[->,thick] (5.1,1.0) -- (5.1,0.15) node[lab,above left=1pt,pos=0.2]{commit};
\fill[red!18] (3.4,0.0) rectangle (5.1,0.15);
\fill[blue!14] (2.9,-0.15) rectangle (3.4,0.0);
\node[lab,red!70!black] at (4.25,0.52)
  {$W_1$: effect durable,\\progress not $\Rightarrow$ duplicate};
\node[lab,blue!60!black] at (1.55,-0.62) {$W_2$: progress\\durable, effect not};
\end{tikzpicture}}
\caption{The EOS boundary on a timeline. Kafka's atomic commit covers the log append
and the offset commit (top); the ledger commit is a separate single-site commit
(bottom). A crash inside either shaded interval yields exactly one durable commit:
$W_1$ duplicates under redelivery, $W_2$ loses the posting permanently. The windows are
defined by Eqs.~\eqref{eq:w1}--\eqref{eq:w2}.}
\label{fig:eos}
\end{figure}

\subsection{Assumptions, crash points, and equivalence classes}
\label{sec:model:assumptions}

Every claim in this paper is conditioned on the assumptions of
Table~\ref{tab:asm}; each is either enforced by the harness or explicitly excluded.
The fault model is crash-recovery, not Byzantine: processes die and restart, storage is
prefix-consistent, and the network may lose, delay, partition or reset connections, but
protocol checksums and retries are intact. Databases and brokers are assumed durable
once they acknowledge; clocks are used only for measurement.

\begin{table*}[t]
\centering
\small
\caption{Assumptions. Every proposition in Section~\ref{sec:formal} is stated under
A1--A13; the enforcement column names the harness mechanism that checks or excludes the
assumption.}
\label{tab:asm}
\begin{tabular}{@{}p{0.7cm}p{10.6cm}p{4.3cm}@{}}
\toprule
ID & Assumption & Enforced / excluded by \\
\midrule
A1 & Crash-recovery faults only; storage writes prefix-consistent; no Byzantine corruption. & process kills and restarts; no data mutation injected \\
A2 & Network faults are loss, delay, partition, reset; protocol checksums and retries intact. & Toxiproxy toxics only \\
A3 & PostgreSQL durable on acknowledgement (\texttt{synchronous\_commit=on}, \texttt{fsync=on}, \texttt{full\_page\_writes=on}); single instance, no replica. & pinned configuration; asserted per run \\
A4 & Kafka durable: replication factor 3, \texttt{min.insync.replicas=2}, \texttt{acks=all}, unclean leader election disabled. & pinned broker configuration \\
A5 & Transactional producer idempotence and fencing as specified; transaction timeout bounded. & pinned producer configuration \\
A6 & At most one owner per partition at a time, except inside the deliberately injected zombie window $C_7$. & pinned consumer timeouts; $C_7$ tests the exception \\
A7 & Fair recovery: killed processes restart and retry until quiescence. & runner restart policy; quiescence criteria \\
A8 & Clocks used only for latency; single host, single boot, no clock steps during a run. & clock checks around every run \\
A9 & Deterministic load; the processor is exactly-once into the source topic. & ground-truth digest and source-integrity gate \\
A10 & Faults act at declared points and every firing is attested. & fault log with one receipt per action \\
A11 & Identical configuration and tuning effort across mechanisms, except declared intrinsic components. & configuration hash audit \\
A12 & Single currency, two-leg entries, positive integer amounts; no reversals or adjustments. & event schema \\
A13 & The oracle shares no code with the pipeline and reads ground truth created before the effect. & independent process and language \\
\bottomrule
\end{tabular}
\end{table*}

Table~\ref{tab:crashpoints} enumerates the crash points of the mechanism skeleton: the
places where a process may die, a message may be lost or duplicated, or two durable
transitions may become visible in an order different from the program order. The list
is exhaustive for the skeleton of Section~\ref{sec:mechanisms} (read, effect, progress,
publish). Points that are outcome-equivalent for the invariants are grouped into the
classes $K_1$--$K_7$ of Table~\ref{tab:kclasses}; classes are used as the cells of the
structural experiment, and the equivalence arguments are the proofs that a class may be
represented by one member.

\begin{table*}[t]
\centering
\small
\caption{Crash-point model. ``Durable at the point'' is the state that survives the
crash; the last column names the window of Eqs.~\eqref{eq:w1}--\eqref{eq:w2}.}
\label{tab:crashpoints}
\begin{tabular}{@{}p{0.6cm}p{6.3cm}p{4.2cm}p{4.6cm}@{}}
\toprule
ID & Crash point & Durable at the point & Window / consequence \\
\midrule
$C_1$ & after fetch, before any durable operation & nothing of the batch & none; replay \\
$C_2$ & inside the effect transaction, before COMMIT is sent & nothing (transaction aborts) & none; replay \\
$C_3$ & after effect COMMIT is sent, before its acknowledgement & effect unknown; progress behind & $W_1$, conditional on the outcome \\
$C_4$ & after effect durable, before progress durable & effect durable; progress behind & $W_1$ (duplicate under redelivery) \\
$C_5$ & after progress durable, before effect durable & progress durable; effect absent & $W_2$ (posting lost permanently) \\
$C_6$ & duplicate delivery without a crash (rebalance reassignment, offset rewind) & unchanged & $W_1$ replay \\
$C_7$ & zombie writer: paused owner resumes after reassignment and writes & the new owner's effect & duplicate write; fencing or dedup needed \\
$C_8$ & publish window: publish sent/lost without a durable publication mark & ledger durable; publication not & publication integrity only \\
$C_9$ & sink-link loss during effect commit (partition, reset) & effect unknown; progress behind & $W_1$, conditional \\
$C_{10}$ & Kafka-side disruption during progress commit (broker kill, coordinator failover, abort) & effect durable; progress not & $W_1$; aborted transaction also discards stream records \\
$C_{11}$ & acknowledgement before durability at either store & unknown & both windows; excluded by A3/A4 \\
$C_{12}$ & partial-batch durability (per-record autocommit) & part of the batch & $W_1$ partial; unreachable with one transaction per batch \\
\bottomrule
\end{tabular}
\end{table*}

\begin{table}[t]
\centering
\small
\caption{Equivalence classes of crash points for the invariants I1--I4. Two points are
equivalent when, for every admissible schedule at one, a schedule exists at the other
with the same durable projection at recovery.}
\label{tab:kclasses}
\begin{tabular}{@{}p{0.7cm}p{1.8cm}p{4.7cm}@{}}
\toprule
Class & Members & Equivalence reason \\
\midrule
$K_1$ & $C_1,C_2$ & no durable trace; identical replay \\
$K_2$ & $C_3,C_9$ & effect possibly durable, progress behind; same unknown-outcome recovery \\
$K_3$ & $C_4,C_{10}$ & effect durable, progress behind; $C_{10}$ additionally aborts stream records \\
$K_4$ & $C_5$ & progress durable, effect absent \\
$K_5$ & $C_6$ & duplicate delivery with progress state ``not yet processed'' \\
$K_6$ & $C_7$ & two concurrent writers of one partition \\
$K_7$ & $C_8$ & ledger and progress durable; publication durability only \\
\bottomrule
\end{tabular}
\end{table}
